\documentclass{article}
\usepackage{amsmath, amssymb, amsthm}
\usepackage{hyperref}
\usepackage{geometry}
\geometry{margin=1in}

\title{Erd\H{o}s Problem \#677}
\date{Last updated: 2025-08-31}
\author{Source: erdosproblems.com}

\begin{document}
\maketitle

\noindent\textbf{Status:} OPEN \quad \textbf{No prize}

\noindent\textbf{Tags:} number theory

\medskip

\noindent\textbf{Problem Statement:}

Let $M(n,k)=[n+1,\ldots,n+k]$ be the least common multiple of $\{n+1,\ldots,n+k\}$. Is it true that for all $m\geq n+k$\[M(n,k) \neq M(m,k)?\]




The Thue-Siegel theorem implies that, for fixed $k$, there are only finitely many $m,n$ such that $m\geq n+k$ and $M(n,k)=M(m,k)$. In general, how many solutions does $M(n,k)=M(m,l)$ have when $m\geq n+k$ and $l&gt;1$? Erd\H{o}s expects very few (and none when $l\geq k$). The only solutions Erd\H{o}s knew were $M(4,3)=M(13,2)$ and $M(3,4)=M(19,2)$. In \cite{Er79d} Erd\H{o}s conjectures the stronger fact that (aside from a finite number of exceptions) if $k&gt;2$ and $m\geq n+k$ then $\prod_{i\leq k}(n+i)$ and $\prod_{i\leq k}(m+i)$ cannot have the same set of prime factors. See also [678] , [686] , and [850] . This is discussed in problem B35 of Guy's collection \cite{Gu04}.




References

[Er79d] Erd\H{o}s, P., Some unconventional problems in number theory . Acta Math. Acad. Sci. Hungar. (1979), 71-80.

[Gu04] Guy, Richard K., Unsolved problems in number theory . (2004), xviii+437.

\noindent\textbf{Related OEIS sequences:} possible


\bigskip
\noindent\small{Source: \url{https://www.erdosproblems.com/677}}
\end{document}
